% !TEX program = xelatex
% Arabic document showing three fonts side by side:
%   Amiri, Noto Kufi Arabic, and Scheherazade New.
% Each paragraph is rendered in its own font, with the font name shown.
\documentclass[12pt, a4paper]{article}

\usepackage{fontspec}
\usepackage{polyglossia}
\setmainlanguage[locale=mashriq]{arabic}
\setotherlanguage{english}

\newfontfamily\arabicfont[Script=Arabic]{Amiri}
\newfontfamily\fontAmiri[Script=Arabic]{Amiri}
\newfontfamily\fontNotoKufi[Script=Arabic]{Noto Kufi Arabic}
\newfontfamily\fontScheherazade[Script=Arabic]{Scheherazade New}

\begin{document}

\title{خطوط عربية في \LR{LaTeX}}
\author{د. نبراس أبو الذهب}
\date{2025}
\maketitle

\section{مقارنة بين الخطوط}

يَستعرضُ هذا المثالُ ثلاثةَ خطوطٍ عربيةٍ شائعةٍ في التنضيدِ الرقميّ.

\subsection{خط Amiri}

{\fontAmiri
خطُ أميري من الخطوطِ النسخيةِ الكلاسيكيةِ المناسبةِ للنصوصِ الأدبيةِ
والقرآنية. صَمّمهُ خالد حسني ليكونَ بديلاً حُرّاً لخطوطِ النسخِ التقليدية.
اسمُ الخطّ: \LR{Amiri}.\par}

\subsection{خط Noto Kufi Arabic}

{\fontNotoKufi
خطُ نوتو الكوفي من خطوطِ شركةِ جوجل، ويتميّزُ بشكلٍ كوفيٍّ عصريٍّ
يناسبُ العناوينَ والواجهات. اسمُ الخطّ: \LR{Noto Kufi Arabic}.\par}

\subsection{خط Scheherazade New}

{\fontScheherazade
خطُ شهرزاد الجديد من خطوطِ مؤسسةِ \LR{SIL}، ويتميّزُ بدعمٍ واسعٍ
للحركاتِ والتشكيلِ ولغاتٍ متعددةٍ تستخدمُ الخطَّ العربيّ.
اسمُ الخطّ: \LR{Scheherazade New}.\par}

\end{document}
